\section{腹膜与原发的类均值乘积}
\label{sec:communication}

这一节把主比较从原发灶点位换成样本。细胞类沿用上一轮的四个程序。乘积用的是类均值，不是把细胞放进CellChat再按细胞数算显著性。7条是看过类均值之后才点名的，所以下面的q不能当成事先锁定的发现。

\begin{figure}[htbp]
  \centering
  \includegraphics[width=\textwidth]{fig_communication.pdf}
  \caption{每份样本一个乘积。蓝点是原发肿瘤，红点是腹膜肿瘤，黑线是原发中位数。横轴是$\log(1+\text{乘积})$。}
  \label{fig:communication}
\end{figure}

C3上皮乘C3AR1巨噬的腹膜中位数是0.542，原发中位数是0.007，3/3高于原发中位数，置换p为0.004，q为0.027（图~\ref{fig:communication}，表~\ref{tab:communication}）。26份里没有一份原发样本超过这个腹膜中位数，但有3份原发超过腹膜里最低的那一份，两组的范围仍有重叠。C3成纤维乘C3AR1巨噬是0.730对0.132，p为0.015，q为0.027；有1份原发样本高于腹膜中位数。两条都用到巨噬C3AR1，不是两个独立发现。

SPP1巨噬乘CD44巨噬是0.505对0.242，只有2/3高于原发中位数，p为0.42，q为0.49。11份原发样本高于腹膜中位数。巨噬SPP1的类均值在腹膜侧较高，但CD44没有一起升高，乘积站不住。SPP1巨噬乘CD44 T细胞是0.964对0.172，3/3，p为0.094，q为0.13，也不过0.05。

CCL2上皮乘CCR2 T细胞是0.021对0.001，3/3，p为0.015，q为0.027。差值只有0.021。CCL2成纤维乘CCR2巨噬是0.009对0.015，方向相反，p为0.80。这条不能写成腹膜里增强的成纤维--巨噬轴。TGFB1 T细胞乘TGFBR2巨噬是0.332对0.136，p为0.014，q为0.027。来源是T细胞，不是成纤维或巨噬。

\begin{table}[htbp]
  \centering
  \caption{样本水平的类均值乘积。p是穷举3份腹膜标签的单侧比例，q是这7条内部的Benjamini--Hochberg。成纤维或上皮缺失的原发样本是sample22，对应组合数3276；其余为3654。}
  \label{tab:communication}
  \footnotesize
  \setlength{\tabcolsep}{3pt}
  \begin{tabular}{lrrrrr}
    \toprule
    乘积 & 腹膜中位数 & 原发中位数 & 高于原发 & p & q \\
    \midrule
    C3成纤维 $\times$ C3AR1巨噬 & 0.730 & 0.132 & 3/3 & 0.015 & 0.027 \\
    C3上皮 $\times$ C3AR1巨噬 & 0.542 & 0.007 & 3/3 & 0.004 & 0.027 \\
    SPP1巨噬 $\times$ CD44巨噬 & 0.505 & 0.242 & 2/3 & 0.420 & 0.490 \\
    SPP1巨噬 $\times$ CD44 T细胞 & 0.964 & 0.172 & 3/3 & 0.094 & 0.131 \\
    CCL2上皮 $\times$ CCR2 T细胞 & 0.021 & 0.001 & 3/3 & 0.015 & 0.027 \\
    CCL2成纤维 $\times$ CCR2巨噬 & 0.009 & 0.015 & 1/3 & 0.797 & 0.797 \\
    TGFB1 T细胞 $\times$ TGFBR2巨噬 & 0.332 & 0.136 & 3/3 & 0.014 & 0.027 \\
    \bottomrule
  \end{tabular}
\end{table}

同一患者的sample38减sample36，C3上皮乘C3AR1巨噬的差是1.177，C3成纤维乘C3AR1巨噬是1.200。这是1对，不另算p。没有把乘积设成0再去预测STAT3。类均值乘积也不是共培养里的阻断。

巨噬STAT3和SOCS3是另一次样本置换，没有并进上面7条的q。STAT3的腹膜中位数是0.641，原发是0.481，3/3，p为0.010；有1份原发高于腹膜中位数。SOCS3是0.695对0.852，方向相反。两者的均值是0.668对0.623，p为0.47。上皮STAT3也是3/3、p为0.015，所以这不是巨噬独有的升高。这是转录本均值，不是磷酸化，也不是把C3--C3AR1删掉之后的下游变化。
